\documentclass[11pt]{amsart}
\usepackage[margin=1in]{geometry}
\usepackage{amssymb,booktabs,hyperref}

\newtheorem{theorem}{Theorem}
\newtheorem{lemma}[theorem]{Lemma}
\newtheorem{proposition}[theorem]{Proposition}
\theoremstyle{remark}
\newtheorem{remark}[theorem]{Remark}

\newcommand{\cp}{\chi_\rho}

\title{The packing chromatic number of the distance graph $D(7,10)$ is between 21 and 29}
\author{AUTHOR NAME}
\date{27 November 2026}

\begin{abstract}
For the integer distance graph $D(7,10)$, we prove
$21\le \cp(D(7,10))\le 29$. The previously published bounds were
$12\le \cp(D(7,10))\le 40$, from Ekstein, Holub and Togni. The upper bound
comes from an explicit periodic packing colouring with period $1088$. The
lower bound is a density argument. Its main ingredient is an integer
potential certifying that colours $1$, $2$ and $4$ together occupy at most
a $67/102$ fraction of any long interval. A periodic partial colouring shows
that this constant is attained. All certificates can be checked with short
standard-library programs. Two independently written implementations
reproduce every number used in the proof.
\end{abstract}

\maketitle

\section{Introduction}

Let $G$ be a graph with distance function $d_G$. A \emph{packing
$k$-colouring} of $G$ is a map $f\colon V(G)\to\{1,\dots,k\}$ such that
$f(u)=f(v)=i$ with $u\ne v$ implies $d_G(u,v)>i$. The \emph{packing
chromatic number} $\cp(G)$ is the least $k$ for which such a colouring
exists. The notion was introduced by Goddard et al.~\cite{GHHHR08} under
the name broadcast chromatic number. Brešar, Ferme, Klavžar and Rall survey
the area in~\cite{BFKR20}.

For positive integers $k<t$, the \emph{distance graph} $D(k,t)$ has vertex
set $\mathbb Z$. Two integers are adjacent when their difference is $\pm k$
or $\pm t$. Packing colourings of distance graphs were studied by
Togni~\cite{Tog14}. For two-element distance sets they were studied by
Ekstein, Holub and Togni~\cite{EHT14}, whose Table~1 lists
$12\le\cp(D(7,10))\le 40$. In that table, the lower bound comes from a
finite subgraph and the upper bound from a periodic construction. Related
work includes $S$-packing colourings of two-step distance
graphs~\cite{BFHJM24}. Closest to the present method is Zhang's use of
finite-state integer certificates for joint colour densities in
$D(1,t)$~\cite{Zha26}. As far as we could find, no improvement of the
$D(7,10)$ entry has been published.

\begin{theorem}\label{thm:main}
$21\le \cp(D(7,10))\le 29$.
\end{theorem}

\section{Distance balls}

Let $B_0=\{0\}$ and
$B_{r+1}=B_r\cup(B_r\pm7)\cup(B_r\pm10)$. By induction, $B_r$ is the set of
integers at graph distance at most $r$ from $0$. Since the graph is
translation invariant, $d(x,y)\le r$ if and only if $y-x\in B_r$.
A colouring is therefore a packing colouring if and only if, for every
colour $c$ and every $d\in B_c\setminus\{0\}$, no two integers at difference
$d$ both receive colour $c$. Every element of $B_r$ has absolute value at most $10r$.

\begin{lemma}\label{lem:ball}
For $8\le c\le 20$ we have $\{1,\dots,10c-15\}\subseteq B_c$, and
$10c-14\notin B_c$.
\end{lemma}
\begin{proof}
Direct finite computation (\texttt{certificates/verify\_lower.cjs}).
\end{proof}

Consequently, two integers of colour $c\ge 8$ are at least $10c-14$ apart, so any interval of length $N$
contains at most $1+\lfloor (N-1)/(10c-14)\rfloor\le N/(10c-14)+1$ of them.

\section{The upper bound}

The file \texttt{certificates/coloring.txt} contains a word
$w=w_0w_1\cdots w_{1087}$ over $\{1,\dots,29\}$, in which every symbol
occurs. Colour $x\in\mathbb Z$ with $w_{x\bmod 1088}$.

\begin{proposition}
This colouring is a packing $29$-colouring of $D(7,10)$.
\end{proposition}
\begin{proof}
Suppose two integers $x<y$ share colour $c$ and $y-x\in B_c$. Translate by a
multiple of $1088$ so that $0\le x<1088$. It therefore suffices to check, for every
residue $i$ with $w_i=c$ and every positive $d\in B_c$, that
$w_{(i+d)\bmod 1088}\ne c$. This finite check passes. It also covers
the case $d\equiv 0\pmod{1088}$, but that case cannot occur because
$\max B_{29}=290<1088$.
\end{proof}

\section{The lower bound}

\subsection{Automata and potentials}

Fix a set $S$ of colours. A \emph{partial $S$-colouring} of an interval
assigns each position a colour in $S$ or a blank $0$, subject to the packing
condition for the colours in $S$. For a single colour $c$, read positions
from left to right and maintain a set $M_c\subseteq\{0,1,\dots,10c-1\}$:
$j\in M_c$ means that some earlier point of colour $c$ forbids colour $c$ at
$j$ positions ahead. Colour $c$ may be placed only if $0\notin M_c$. To advance,
subtract $1$ from every element and discard negatives. If $c$ was placed,
add $\{d-1: d\in B_c,\ d>0\}$. For several colours, run one such
automaton per colour, and allow at most one placement per step. Starting
from the empty state, the words accepted are exactly the partial
$S$-colourings of an interval. Restricting a packing colouring of $\mathbb Z$ to an interval $I$, and blanking the colours outside $S$,
gives such a word of length $|I|$.

Give a transition weight $w=1$ if it places a colour and $w=0$ otherwise.

\begin{lemma}\label{lem:potential}
Let $\mathcal S$ be a set of states that contains the empty state $s_0$ and
is closed under all legal transitions. Let $p,q$ be positive integers. Suppose
$H\colon\mathcal S\to\mathbb Z_{\ge0}$ satisfies
\[ q\,w(s,t)-p\le H(s)-H(t) \]
for every legal transition $s\to t$. Then every partial $S$-colouring of an
interval $I$ has at most $\frac pq|I|+\frac{H(s_0)}{q}$ coloured positions.
\end{lemma}
\begin{proof}
Sum the inequality along the run $s_0\to s_1\to\dots\to s_{|I|}$. The
right-hand side telescopes to $H(s_0)-H(s_{|I|})\le H(s_0)$.
\end{proof}

All bounds below have the form $N_S(I)\le r_S|I|+C_S$ with $C_S$
independent of $I$ and of the colouring. Here $N_S(I)$ is the number of positions of $I$ with a colour in $S$. They can therefore be added on a common interval, and no
assumption that colour densities exist is needed.

\subsection{Colours 1, 2 and 4 together}

\begin{proposition}\label{prop:124}
In every packing colouring of $D(7,10)$ and every interval $I$,
\[ 102\,N_{\{1,2,4\}}(I)\le 67|I|+569. \]
The constant $67/102$ is best possible.
\end{proposition}
\begin{proof}
The single-colour automata for colours $1$, $2$ and $4$ have $169$, $639$
and $533$ states. Their product, restricted to the states reachable from the empty
state, has $24\,274\,996$ states and $45\,060\,505$ legal transitions. The
file \texttt{certificates/potential-124.u16le.gz} assigns a value in
$\{0,\dots,65534\}$ to each reachable state, and $65535$ to the others. The
verifier rebuilds the automata from the step set $\{\pm7,\pm10\}$. It checks
that $s_0$ has a value and that every legal successor of a valued state has a value,
and it checks the inequality of Lemma~\ref{lem:potential} with $(p,q)=(67,102)$
on every transition. The value at $s_0$ is $569$.

For optimality, \texttt{certificates/density-124.txt} is a word of length
$102$ over $\{0,1,2,4\}$ with $67$ non-blank letters. Repeated periodically,
it passes the same check as in the upper bound.
\end{proof}

\begin{remark}
Taken separately, the exact maximum densities of colours $1$, $2$, $4$ are
$8/17$, $7/37$ and $3/41$ (computed by Karp's maximum mean cycle
algorithm~\cite{Kar78}). Their sum is about $0.733$, while
$67/102\approx0.657$. The joint certificate is therefore essential.
\end{remark}

\subsection{Colours 3, 5, 6, 7}

For each of these colours the single-colour automaton is small. Iterating
$H_{j+1}(s)=\max\bigl(H_j(s),\max_{s\to t}\{H_j(t)+q\,w(s,t)-p\}\bigr)$ from
$H_0\equiv0$ reaches a fixed point. A fixed point satisfies the hypothesis of
Lemma~\ref{lem:potential}.

\begin{center}
\begin{tabular}{cccc}
\toprule
colour & states & $p/q$ & $H(s_0)$\\
\midrule
3 & 558 & $2/17$ & 45\\
5 & 273 & $2/45$ & 84\\
6 & 289 & $2/55$ & 104\\
7 & 289 & $1/34$ & 62\\
\bottomrule
\end{tabular}
\end{center}

Each of these four ratios equals the maximum cycle mean of the
corresponding automaton, so they are the exact maximum densities.

\subsection{Conclusion of the proof}

Suppose $f$ is a packing $20$-colouring of $D(7,10)$, and let $I$ be an
interval of length $N$. Every position carries one of the colours
$1,\dots,20$. Adding Proposition~\ref{prop:124}, the table above and
Lemma~\ref{lem:ball} gives
\[ N\le \sigma N + C, \]
where $C$ is a constant independent of $N$ and
\[
\sigma=\frac{67}{102}+\frac2{17}+\frac2{45}+\frac2{55}+\frac1{34}
+\sum_{c=8}^{20}\frac1{10c-14}
=\frac{128279378457743209}{128437412456857504}<1 .
\]
For large $N$ this is impossible. Hence $\cp(D(7,10))\ge21$, which completes
the proof of Theorem~\ref{thm:main}. \qed

\begin{remark}
The margin is small: $1-\sigma\approx1.23\times10^{-3}$. Replacing
$67/102$ by the sum of the separate maxima of colours $1,2,4$ would give
$\sigma>1$. Joint certificates for larger colour sets may raise the lower
bound further.
\end{remark}

\section{Reproducibility}

The repository accompanying this note contains:
\begin{itemize}
\item \texttt{certificates/}: the two periodic words, the compressed potential, and the
verifiers \texttt{verify\_upper.py} (Python~3) and \texttt{verify\_lower.cjs}
(Node.js). Both use only the standard library. The lower-bound
verifier runs in about a second.
\item \texttt{independent-check/}: a second implementation, written
separately. It checks the periodic words with breadth-first graph
distances instead of balls. It recomputes the single-colour maximum
densities with Karp's algorithm. It also recomputes the joint
$\{1,2,4\}$ potential from scratch by value iteration, without using the
supplied certificate, and obtains the same state count, transition count and
$H(s_0)=569$.
\end{itemize}

\section*{Acknowledgements}

The search for the colouring and the certificates was carried out with AI
coding agents (OpenAI Codex/GPT). The independent re-verification was done
with Anthropic's Claude. The author checked the results and is responsible for
the content.

\begin{thebibliography}{9}
\bibitem{BFHJM24} B.~Brešar, J.~Ferme, P.~Holub, M.~Jakovac, P.~Melicharová,
$S$-packing colorings of distance graphs with distance sets of cardinality
2, arXiv:2405.18904, 2024.
\bibitem{BFKR20} B.~Brešar, J.~Ferme, S.~Klavžar, D.~F.~Rall, A survey on
packing colorings, \emph{Discuss. Math. Graph Theory} 40 (2020) 923--970,
doi:10.7151/dmgt.2320.
\bibitem{EHT14} J.~Ekstein, P.~Holub, O.~Togni, The packing coloring of
distance graphs $D(k,t)$, \emph{Discrete Appl. Math.} 167 (2014) 100--106,
doi:10.1016/j.dam.2013.10.036; arXiv:1302.0721.
\bibitem{GHHHR08} W.~Goddard, S.~M.~Hedetniemi, S.~T.~Hedetniemi,
J.~M.~Harris, D.~F.~Rall, Broadcast chromatic numbers of graphs,
\emph{Ars Combin.} 86 (2008) 33--49.
\bibitem{Kar78} R.~M.~Karp, A characterization of the minimum cycle mean in
a digraph, \emph{Discrete Math.} 23 (1978) 309--311.
\bibitem{Tog14} O.~Togni, On packing colorings of distance graphs,
\emph{Discrete Appl. Math.} 167 (2014) 280--289.
\bibitem{Zha26} E.~Zhang, Density regions, integer certificates and packing
colorings of distance graphs, arXiv:2609.09018, 2026.
\end{thebibliography}

\end{document}
